% Pillar-7: E1 open-backward-pass gradient validation (toy scale, directional)
\section{E1: An Open-Backward-Pass Check of T3}
\label{sec:p7-e1}
Theorems T1 and T2 are checkable from reward telemetry alone, but T3---that
gradient magnitude tracks $p(1-p)$---requires seeing the gradients, which
the closed stack used for most of the benchmark does not expose. Experiment
E1 (W\&B project \texttt{zvf-colab-experiments}, run family
\texttt{E1\_grad\_signal}) therefore reruns the loop in an open trainer with
a fully instrumented backward pass: a 0.5B-parameter model on a synthetic
arithmetic task with binary verifiable rewards, chosen so that per-step
success probability sweeps a wide range of $p$ within a single run.

At each step we log the empirical group success rate $\hat p$, the contrast
term $\hat p(1-\hat p)$, and the global gradient norm $\|\nabla\|$ after the
GRPO loss. The result:
\begin{equation}
  \mathrm{corr}\bigl(\|\nabla\|,\; \hat p(1-\hat p)\bigr) \;=\; +0.71 .
  \label{eq:p7-e1}
\end{equation}
The sign and strength are what T3 predicts: steps whose groups carry more
within-group contrast produce proportionally larger policy-gradient norms,
and steps dominated by degenerate groups produce gradients near zero even
when the raw batch reward is high.

\paragraph{Honest scope.} We state plainly what E1 is and is not. It is a
\emph{directional} validation on a 0.5B toy model and a synthetic arithmetic
task---two ingredients chosen precisely to make the backward pass cheap and
inspectable, and both far from the production regime. The $+0.71$ is a
single-configuration correlation, not a publishable effect size: it comes
with no confidence interval across model scales, no ablation over loss
variants, and no evidence that the coefficient (as opposed to the sign)
transfers to 8B--32B models or natural-language reasoning data. We use E1
for exactly one inferential step: it makes the $p(1-p)$ factor in
Eq.~\eqref{eq:p7-S} an observed property of at least one real backward pass
rather than a modeling convenience. The decisive version of this
experiment---per-group gradient attribution at production scale, across
seeds---is specified in Section~\ref{sec:p7-limitations}.
